% Part: history
% Chapter: biographies 
% Section: alfred-tarski

\documentclass[../../../include/open-logic-section]{subfiles}

\begin{document}

\olfileid[hi]{his}{bio}{tar}

\olsection{अल्फ्रेड टार्स्की}

\olphoto{tarski-alfred}{अल्फ्रेड टार्स्की}

अल्फ्रेड टार्स्की का जन्म 14 जनवरी 1901 को पोलैंड के वारसॉ (तब रूसी
साम्राज्य का भाग) में हुआ। ``नेपोलियन-जैसे'' कहे जाने वाले टार्स्की ऊँची
आवाज़ वाले, वाचाल और प्रखर थे। उनकी ऊर्जा प्रायः उनके व्याख्यानों में दिखती
थी---एक बार व्याख्यान के दौरान सिगरेट फेंकते हुए उन्होंने कूड़ेदान में आग
लगा दी, जिसके बाद उस भवन में उनके व्याख्यान देने पर रोक लगा दी गई।

टार्स्की में कम आयु से ही ज्ञान की प्यास थी। यद्यपि बाद के जीवन में वे
विद्यार्थियों को कहते थे कि उन्होंने तर्कशास्त्र इसलिए पढ़ा क्योंकि वही एक
कक्षा थी जिसमें उन्हें~B मिला था, किंतु उनके माध्यमिक विद्यालय के अभिलेख
दिखाते हैं कि उन्हें हर विषय में A मिला था---तर्कशास्त्र में भी। उन्होंने
1918 से 1924 तक वारसॉ विश्वविद्यालय में अध्ययन किया। टार्स्की पहले जीवविज्ञान
पढ़ना चाहते थे, किंतु विश्वविद्यालय तर्क और दर्शन के वारसॉ संप्रदाय का
केंद्र होने के कारण उनकी रुचि गणित, दर्शन और तर्कशास्त्र में हो गई। टार्स्की
ने 1924 में स्तानिस्वाव लेश्न्येव्स्की के निर्देशन में डॉक्टरेट प्राप्त की।

1939 में संयुक्त राज्य अमेरिका जाने से पहले टार्स्की ने वारसॉ में माध्यमिक
विद्यालय के शिक्षक के रूप में काम करते हुए अपने कुछ सबसे महत्त्वपूर्ण कार्य
पूरे किए। तार्किक परिणाम और तार्किक सत्य पर उनका कार्य इसी समय लिखा गया।
1939 में टार्स्की व्याख्यान यात्रा पर संयुक्त राज्य अमेरिका गए हुए थे। उसी
दौरान जर्मनी ने पोलैंड पर आक्रमण किया और यहूदी मूल के कारण टार्स्की लौट नहीं
सके। उनकी पत्नी और बच्चे युद्ध के अंत तक पोलैंड में रहे, फिर वे भी संयुक्त
राज्य अमेरिका आ सके। टार्स्की ने हार्वर्ड, कॉलेज ऑफ द सिटी ऑफ न्यू यॉर्क,
प्रिंसटन के इंस्टीट्यूट फॉर एडवांस्ड स्टडी और अंततः यूनिवर्सिटी ऑफ
कैलिफोर्निया, बर्कली में पढ़ाया। वहाँ उन्होंने तर्क और विज्ञान की पद्धति का
बहुविषयी कार्यक्रम स्थापित किया। टार्स्की का निधन 26 अक्टूबर 1983 को 82 वर्ष
की आयु में हुआ।

\begin{reading} 
टार्स्की के जीवन के बारे में अधिक जानने के लिए जीवनी \emph{Alfred Tarski:
Life and Logic} \citep{Feferman2004} देखें। तार्किक परिणाम और सत्य पर
टार्स्की की मूलभूत रचनाएँ \citep{Tarski1983} में अंग्रेज़ी में उपलब्ध हैं।
टार्स्की की सभी मूल रचनाएँ चार-खंडों की शृंखला \citep{Tarski1981} में
संकलित की गई हैं।
\end{reading}

\end{document}
